\section{Approximation consequences}\label{sec:consequences}

Theorem~\ref{thm:ugc} turns established conditional hardness results
into NP-hardness results. The underlying conditional reductions and
approximation thresholds are due to the cited authors. We give the
additional arguments needed for the exact instance models stated below.
We also record two results from companion
manuscripts. For a maximization problem, ratio
\(r\) means returning value at least \(r\) times the optimum; for a
minimization problem, factor \(c\) means returning value at most \(c\)
times the optimum.
All factors and tolerances below are fixed independently of input size.
In each reduction using Theorem~\ref{thm:ugc}, they determine the
Unique Games errors before its alphabet is chosen. If
\(\mathrm P\ne\mathrm{NP}\), the resulting hardness statements exclude
deterministic polynomial-time approximations at the indicated strict
thresholds.

\subsection{Constraint satisfaction and Vertex Cover}

For \textsc{Max-Cut}, which maximizes the number of edges crossing a
bipartition, Khot, Kindler, Mossel, and O'Donnell's reduction gives
NP-hardness of approximation at every fixed ratio in
\((\alpha_{\mathrm{GW}},1)\), where the Goemans--Williamson constant is
\(\alpha_{\mathrm{GW}}\simeq0.87856\)~\cite[Theorem~1]{KKMO07}.
This matches the semidefinite approximation threshold of Goemans and
Williamson~\cite{GW95}. The reduction uses the Majority Is Stablest
theorem of Mossel, O'Donnell, and Oleszkiewicz~\cite{MOO10}.
For \textsc{Vertex Cover}, which minimizes the number of vertices
meeting every edge, Khot and Regev's reduction gives NP-hardness at
every fixed factor in \((1,2)\)~\cite{KhotRegev2008}.

A constraint satisfaction problem (CSP) is specified here by a finite
nonempty domain \(D\) and a finite family \(\Lambda\) of Boolean
predicates on \(D\), of bounded arity. A weighted \(\Lambda\)-instance
has a finite variable set and a finite collection of predicates from
\(\Lambda\) applied to tuples of variables, with nonnegative rational
weights summing to one. An assignment maps the variables to \(D\);
its value is the total satisfied weight, and \(\operatorname{OPT}\)
is the maximum of this value over all assignments.
Raghavendra's theorem characterizes the optimal approximation
for every fixed language of this form, up to arbitrarily small positive
error, by a basic semidefinite relaxation
\cite[Theorem~1.1 and Corollary~1.3]{Raghavendra08}.
The Vertex Cover result uses its own reduction. Write
\(\operatorname{SDP}_{R}\) for the value of Raghavendra's relaxation,
as defined in~\cite[Section~4.5]{RaghavendraThesis09}.

The fixed-instance form of Raghavendra's theorem and the ordering
theorem below use a strengthened bipartite formulation of Unique Games.
In its YES case, one labeling satisfies every edge incident to at least
a \(1-\zeta\) fraction of left vertices; in its NO case, every labeling
satisfies at most a \(\nu\) fraction of edges. These are uniform
fractions of the left vertices and edges, respectively.
Lemma~\ref{lem:simple-strong-left} gives this promise on an explicit
nonempty simple graph whose common positive left degree can meet any
fixed lower bound and be divisible by any fixed positive integer. It follows
Khot and Regev's normalization and sequence construction
\cite[Lemmas~3.3, 3.4, and~3.6]{KhotRegev2008}, with an additional
restriction that makes the graph simple. Its error parameters and
degree are fixed before the alphabet is chosen, so
Theorem~\ref{thm:ugc} supplies the required premise.

\begin{corollary}[Raghavendra's consequence]\label{cor:raghavendra-consequence}
Let \(J\) be a fixed \(\Lambda\)-instance with
\(A=\operatorname{OPT}(J)<B=\operatorname{SDP}_{R}(J)\).
For every fixed \(0<\eta<(B-A)/2\), it is NP-hard to distinguish
\(\Lambda\)-instances \(I\) satisfying
\[
 \operatorname{OPT}(I)\ge B-\eta
 \qquad\text{from those satisfying}\qquad
 \operatorname{OPT}(I)\le A+\eta.
\]
\end{corollary}
\begin{proof}
For the thesis's common-arity convention
\cite[Definitions~7.2.1--7.2.2]{RaghavendraThesis09}, identify \(D\)
with a fixed \([q]\), and pad every predicate with ignored trailing
arguments to a common positive arity, filling the added positions with
a fixed variable of \(J\), which exists because \(A<B\). This leaves
each applied payoff and its dependence set unchanged, hence preserves
\(\operatorname{OPT}\) and the relaxation of
\cite[Section~4.5]{RaghavendraThesis09}.
In that relaxation, each local distribution is over assignments to
the variables on which its applied payoff actually depends.
When emitting an occurrence, discard only the added trailing positions,
retaining the original predicate from \(\Lambda\) and its weight.

Apply Raghavendra's fixed-instance theorem
\cite[Theorem~7.9]{RaghavendraThesis09} with tolerance \(\eta/2\);
see also \cite[Theorem~1.1 and Lemma~2.3]{Raghavendra08}.
The fixed witness and tolerance determine its strengthened Unique Games
errors, so the formulation above and Theorem~\ref{thm:ugc} supply its
premise, with a fixed alphabet \([M]\). Its verifier instances have
optimum greater than \(B-\eta/2\) in the YES case and at most
\(A+\eta/2\) in the NO case.

We pass from the verifier's probability weights to the stated rational
weights. For this fixed invocation, the construction in
\cite[Sections~7.3.1 and~7.5]{RaghavendraThesis09} gives a local
probability law \(\theta\) on a finite set \(\mathcal S\), after the
alphabet \([M]\) is fixed. A record \(s\in\mathcal S\) specifies an
applied predicate of \(J\), its \(a(s)\) ordered oracle-query strings
in \(D^M\), and the positions filled by their replies. This law may
depend on \(J,\eta,M\), but no varying input parameter enters it.
Rational probability vectors are dense in this finite simplex, so fix
\(\widehat\theta\in\mathbb Q_{\ge0}^{\mathcal S}\) with
\[
 \sum_{s\in\mathcal S}\widehat\theta(s)=1,
 \qquad
 \delta:=\frac12\sum_{s\in\mathcal S}
       |\widehat\theta(s)-\theta(s)|<\frac{\eta}{2}.
\]
This rational table is fixed finite data of the reduction; its
description length may depend on \(J,\eta,M\).

Let \(L_\Phi\) and \(V_\Phi\) be the left and right sides of the
nonempty simple game supplied above, and let \(h\ge1\) be its common
left degree.
List the variables \(V_\Phi\times D^M\) used in Section~7.5 of the
cited thesis. Its outer verifier chooses a uniform left vertex \(u\)
and, independently for each
oracle query, chooses a uniform neighbor of \(u\) and relabels the
query string by that edge's permutation. For each \(s\in\mathcal S\),
\(u\in L_\Phi\), and ordered neighbor tuple in \(N(u)^{a(s)}\), emit
the \(\Lambda\)-predicate prescribed by \(s\) on the resulting query
variables. Fill all positions of any syntactic variable omitted from
the test with the first listed output variable; the applied predicate
is independent of such variables, so this leaves its value unchanged.
Give the occurrence
weight
\[
 \frac{\widehat\theta(s)}{|L_\Phi|h^{a(s)}}.
\]
These weights are nonnegative rationals summing to one. Keeping equal
occurrences separate gives
\(|L_\Phi|\sum_{s\in\mathcal S}h^{a(s)}\) occurrences and
\(|V_\Phi||D|^M\) variables. The local set and every \(a(s)\) are
fixed, while the graph has polynomial size; hence this is a polynomial
occurrence list. The explicit neighbor lists and edge permutations let
us write it in polynomial time. Each weight has polynomial binary
length: it can be represented with a fixed denominator multiplied by
\(|L_\Phi|h^{a(s)}\).

For an assignment \(\sigma\), let \(g_\sigma(s)\in[0,1]\) be its
conditional acceptance probability over these outer choices. The
original and rationalized instances use the same variable set, and
\[
 \bigl|\operatorname{val}_{\widehat\theta}(\sigma)
       -\operatorname{val}_{\theta}(\sigma)\bigr|
 =\left|\sum_{s\in\mathcal S}
       \bigl(\widehat\theta(s)-\theta(s)\bigr)g_\sigma(s)\right|
 \le\delta<\frac{\eta}{2}.
\]
Taking maxima gives the same bound for their optima. The imported gap
therefore yields optimum at least \(B-\eta\) in the YES case and at most
\(A+\eta\) in the NO case, as required.
\end{proof}

\subsection{Ordering constraints}

An ordering CSP uses a global order of the input variables, rather than
an assignment from a fixed finite domain. Fix an integer \(k\ge2\) and
a nonempty proper subset \(\Pi\subset S_k\), where \(S_k\) is the set
of permutations of \(\{1,\ldots,k\}\). In the finite
probability-weighted model of Guruswami, H\aa stad, Manokaran,
Raghavendra, and Charikar~\cite[Definition~8.2]{GuruswamiEtAl2011OrderingCSP},
we use an explicitly listed finite variable set \(V\) and a finite list of
constraint occurrences. Their nonnegative rational weights are encoded
by binary numerators and denominators and sum to one. Every occurrence is an
ordered \(k\)-tuple \((v_1,\ldots,v_k)\) of distinct variables,
satisfied by a linear order \(\prec\) when
\[
 v_{\pi(1)}\prec\cdots\prec v_{\pi(k)}
 \quad\text{for some }\pi\in\Pi.
\]
Write \(\OPT_\Pi(I)\) for the maximum satisfied weight. A uniform
random order satisfies expected weight \(\rho_\Pi=|\Pi|/k!\).
Since \(\OPT_\Pi(I)\le1\), it is a \(\rho_\Pi\)-approximation.
For fixed \(k\), conditional expectation derandomizes this guarantee.
After any prefix of the order has been exposed, a constraint has at
most \(k!\) possible relative orders consistent with that prefix.
Its conditional success probability is therefore computable exactly;
choosing each next variable to preserve the expected value takes a
polynomial number of exact rational operations of polynomial bit length.
The size of \(V\) is part
of the input, so this is not a CSP over the fixed finite domain used
in the preceding corollary.

\begin{corollary}[Ordering-CSP consequence]\label{cor:ordering-csp}
Fix \(k\ge2\) and \(\varnothing\ne\Pi\subsetneq S_k\), and put
\(\rho_\Pi=|\Pi|/k!\). For every fixed
\(0<\varepsilon<(1-\rho_\Pi)/2\), it is NP-hard to distinguish finite
weighted \(\Pi\)-ordering-CSP instances \(I\) satisfying
\[
 \OPT_\Pi(I)\ge1-\varepsilon
 \qquad\text{from those satisfying}\qquad
 \OPT_\Pi(I)\le\rho_\Pi+\varepsilon.
\]
Consequently, approximation at every fixed ratio
\(r\in(\rho_\Pi,1]\) is NP-hard. Together with the random-ordering
guarantee, this gives the optimal threshold \(\rho_\Pi\).
\end{corollary}
\begin{proof}
Proposition~\ref{prop:finite-ordering-transfer} constructs the stated
finite instance from a simple strong-left Unique Games instance. It
uses the local gap and analytic estimates of
Guruswami et al.~\cite{GuruswamiEtAl2011OrderingCSP} and fixes the
required game errors and degree conditions before the alphabet is
chosen. Lemma~\ref{lem:simple-strong-left} supplies that game from
Theorem~\ref{thm:ugc}; composing the reductions proves the displayed
gap for the original predicate \(\Pi\).

For a fixed \(r\in(\rho_\Pi,1]\), choose the positive error still
smaller so that \(\varepsilon<(r-\rho_\Pi)/(1+r)\). Then
\(r(1-\varepsilon)>\rho_\Pi+\varepsilon\), so an
\(r\)-approximation would distinguish the two cases.
\end{proof}

For \textsc{Maximum Acyclic Subgraph} in the finite weighted
directed-graph model, take \(k=2\) and \(\Pi=\{12\}\) in one-line
notation. A constraint \((u,v)\) asks that the edge point forward;
forward edges form an acyclic subgraph, and every acyclic subgraph has
a topological order. Its threshold is therefore \(1/2\).
For \textsc{Betweenness}, the constraint \((a,b,c)\) asks that \(b\)
lie between \(a\) and \(c\), so \(\Pi=\{123,321\}\subset S_3\) and
the threshold is \(2/3!=1/3\)
\cite[Theorem~1.1 and Section~1.2]{GuruswamiEtAl2011OrderingCSP}.
These finite weighted consequences have near-perfect completeness
with fixed positive error.

\subsection{Cut, deletion, and clustering consequences}

For the cut problems of Chawla, Krauthgamer, Kumar, Rabani, and
Sivakumar~\cite[Section~1]{CKKRS2006}, let \(G=(V,E)\) be a finite
undirected graph with positive rational edge costs \(c_e\), and let
\(\mathcal D=\{\{s_i,t_i\}:1\le i\le k\}\) contain \(k\ge1\)
distinct demand pairs with \(s_i\ne t_i\). For \(M\subseteq E\), write
\(c(M)=\sum_{e\in M}c_e\), and let \(q_{\mathcal D}(M)\) count
the demand pairs disconnected in \(G-M\).
Weighted \textsc{Multicut} minimizes \(c(M)\) subject to
\(q_{\mathcal D}(M)=k\). Their nonuniform \textsc{Sparsest Cut}
search problem minimizes \(c(M)/q_{\mathcal D}(M)\) over cutsets
\(M\) with \(q_{\mathcal D}(M)>0\). Only the selected pairs carry
unit demand; they need not be all vertex pairs.

For weighted Min-\(2\mathrm{CNF}^{\equiv}\) Deletion, the input
consists of clauses \((\ell\equiv\ell')\) with positive rational
weights, where \(\ell,\ell'\) are Boolean literals, and an assignment
pays the total weight of false clauses. The allowed clauses therefore
express equality or disequality of variables.

In the arbitrary weighted-graph form of \textsc{Correlation Clustering},
each present edge has a \(+\) or \(-\) label and a nonnegative rational
weight. A partition into any number of clusters pays for \(+\) edges
whose endpoints lie in different clusters and \(-\) edges whose
endpoints lie in the same cluster. Missing edges carry no penalty.
This is the minimum-disagreements objective of
\cite[Section~2.1]{DemaineEtAl2006Clustering}.

\begin{corollary}[Cut, deletion, and clustering hardness]
\label{cor:cut-clustering}
For every fixed real constant \(C>1\), it is NP-hard to approximate
each of the four minimization problems above within factor \(C\).
For \textsc{Sparsest Cut}, this is hardness of the search task under
polynomial-time Cook reductions: the algorithm must output a cutset,
not only an estimate of the optimum. Consequently, unless
\(\mathrm P=\mathrm{NP}\), none of these problems admits a deterministic
polynomial-time approximation with any fixed constant factor.
\end{corollary}
\begin{proof}
Corollary~1.3 of \cite{CKKRS2006} gives the first three conclusions
under its ordinary edge-value Unique Games hypothesis.
For a requested factor \(C\), use its hardness statement at any fixed
larger factor \(L\): a \(C\)-approximation would also be an
\(L\)-approximation.
To supply that hypothesis, orient the bipartite constraints of
Theorem~\ref{thm:ugc} left-to-right, renumber its fixed alphabet as
\([d]\) with \(d=2^s\ge2\), and give its edges equal weights summing
to one. Padding the smaller vertex side with isolated questions if
needed preserves both value and polynomial size. The search-version
qualification for \textsc{Sparsest Cut} is in Section~1.3 of the
cited source.

For \textsc{Correlation Clustering}, the hard \textsc{Multicut}
instances in~\cite[Section~2.1]{CKKRS2006} have distinct antipodal
demand pairs inside \(d\)-cubes, while the edges of the Multicut graph
are cube edges or join different cubes. Here \(d\ge2\), unchanged by
Lemma~1.5 of
that source, so the demand pairs are nonedges. The transformation of
Demaine, Emanuel, Fiat, and Immorlica
\cite[Theorem~4.7 and Corollary~4.8]{DemaineEtAl2006Clustering}
therefore creates no opposite-signed parallel edges. It preserves the
optimum and converts every clustering into a multicut of no greater
cost, so the same fixed-factor hardness transfers.
\end{proof}

The separate manuscript
\cite[Theorem~1.1]{OpenAIUniformSparsest2026} proves
NP-hardness at every fixed factor \(C>1\) for the bipartition objective with
nonnegative rational capacities and exactly unit demand between every
pair of distinct vertices. This uniform-demand theorem does not depend on
Theorem~\ref{thm:ugc} and is not supplied by the nonuniform reduction
above.

\paragraph{Kernel clustering.}
A separate companion application concerns the following maximization
problem. For \emph{identity-target kernel clustering}, fix an integer \(k\ge3\).
An input is an explicitly represented rational symmetric matrix
\(A=(a_{pq})_{p,q=1}^N\succeq0\) with \(A\mathbf1=0\), where
\(\mathbf1\) is the all-ones vector.
An assignment \(\sigma:\{1,\ldots,N\}\to\{1,\ldots,k\}\) has value
\[
 \val_A(\sigma)=\sum_{p,q=1}^N
 a_{pq}\mathbf1_{\{\sigma(p)=\sigma(q)\}},
 \qquad
 \OPT_k(A)=\max_\sigma\val_A(\sigma).
\]
The sum includes ordered pairs and the diagonal; empty clusters are
allowed, and there is no normalization by cluster size. A loss factor
\(\alpha\ge1\) means returning value at least \(\OPT_k(A)/\alpha\).
The Gaussian propeller companion uses Theorem~\ref{thm:ugc} to prove
NP-hardness of achieving every fixed loss factor
\[
 1\le\alpha<\alpha_k,
 \qquad \alpha_k=\frac{8\pi}{9}\left(1-\frac1k\right),
\]
for each fixed \(k\ge3\)
\cite[Theorem~6.1]{OpenAIGaussianPropeller2026}.
This is strict hardness below the propeller rounding factor
\(\alpha_k\).
